    The Cooley-Tukey FFT works when the $n$ points to sample are $n$-th roots of unity, and hence form a cyclic multiplicative group;
    this and similar algorithms are also referred to as \textit{multiplicative} FFTs, as they are based on the factorization of $x^n-1$.\\
    Unfortunately, multiplicative FFTs do not work when the underlying field does not admit primitive roots of unity, as it happens for fields
    of characteristic two; for this reason, \textit{additive} FFTs were invented~\cite{wang1988fast,cantor1989arithmetical}.
    Instead of using the $n$-th roots of unity, additive FFTs evaluate the polynomials at the powers of a primitive $(n-1)$-th root of unity,
    as well as the zero element. Some examples of such algorithms that work on binary fields are the Gao-Mateer FFT~\cite{gao2010additive} 
    and the Lin-Chung-Han FFT~\cite{lin2014novel}.\\
    We will now focus on the LCH FFT, as it will be the foundation for a faster approach based on the Frobenius map; for this reason, we will
    consider the case $n=2^m$ where $m$ is a power of 2.


\subsection{Cantor basis}
In~\cite{cantor1989arithmetical} Cantor showed how to get a special basis in which it is possible to compute the FFT efficiently
    for any $m = p^n$, with $p$ a prime number; later in~\cite{gao2010additive} Gao and Mateer used a simpler construction for the special case
    $p=2$, where the Cantor basis $\mathbf{v} = (v_1, \cdots, v_{m})$ of $\mathds{F}_2^m$ follows the structure:
    \begin{equation}
        v_1 = 1, \quad v_i^2+v_i = v_{i-1} \text{ for } i > 1
    \end{equation}
    With respect to the Cantor basis $\mathbf{v}$, we define its sequence of subspaces as:
    \begin{equation}
        W_0 := \{0\}, \quad W_i := span\{v_0, \cdots, v_{i-1}\} \text{ for } i \geqslant 1
    \end{equation}

    By construction, we have that $W_0 \subset W_1 \subset \cdots \subset W_{m-1} = \mathds{F}_{2^m}$ 
    and the basis for a smaller space is contained into the basis of a larger space: this means that 
    we can arbitrarily treat elements of smaller fields as elements in the larger fields by padding zeros to the extra dimensions.

    With respect to each element of the basis $v_i$ we can define its subspace polynomial $s_i(x)$, also called vanishing polynomial, as:
    \begin{equation}
        s_i(x) = \prod_{a \in W_i} (x - a)
    \end{equation}
    Since $\dim(W_i) = i$ we have $\deg(s_i(x)) = |W_i| = 2^i$.
    Such polynomials, as shown in~\cite{cantor1989arithmetical,gao2010additive}, satisfy the following useful properties:
    \begin{itemize}
        \item linearity: $s_i(x+y) = s_i(x) + s_i(y)$
        \item recursivity: $s_{i+j}(x) = s_i(s_j(x))$ or, equivalently, $s_{i+1}(x) = s_i(x)^2+s_i(x)$
        \item $s_i(v_i) = 1$, for $v_i$ in the Cantor basis
        \item the roots of $s_i(x)$ are the elements of $W_i$
    \end{itemize}
    Note that the $W_i$s are also fields whenever $i$ is a power of two; in that case, the vanishing polynomial has the form $s_i(x) = x^{2^i}+x$.

    \begin{example}[Cantor basis in $\mathds{F}_{16}$]
        Recalling that $GF(16) = GF(2)[X]/\langle X^4+X+1 \rangle$ and that $v_1 = 1$, we can compute $v_i = \beta(1) + \gamma(x) + \delta(x^2) + \epsilon(x^3)$,
        where the coefficients can be found by solving:
        \begin{equation*}
            \beta(1+1^2)+\gamma(x+x^2) + \delta(x^2+(x^2)^2) + \epsilon(x^3+(x^3)^2) = v_{i-1}
        \end{equation*}
        The resulting Cantor basis is $\mathbf{v} = (1, \alpha+\alpha^2, \alpha, \alpha+\alpha^3)$. As a consequence, the subspaces $W_i$ are:
        \begin{itemize}
            \item $W_0 = \{0\}$
            \item $W_1 = \{0, 1\}$
            \item $W_2 = \{0, 1, \alpha+\alpha^2, 1+\alpha+\alpha^2\}$
            \item $W_3 = \{0, 1, \alpha+\alpha^2, 1+\alpha+\alpha^2, \alpha, 1+\alpha, \alpha^2, 1+\alpha^2\}$
            \item $W_4 = \{0, 1, \alpha+\alpha^2, 1+\alpha+\alpha^2, \alpha, 1+\alpha, \alpha^2, 1+\alpha^2, \alpha+\alpha^3, 1+\alpha+\alpha^3, \alpha^2+\alpha^3, 1+\alpha^2+\alpha^3, \alpha^3, 1+\alpha^3, \alpha+\alpha^2+\alpha^3, 1+\alpha+\alpha^2+\alpha^3\} = \mathds{F}_{16}$
        \end{itemize}
    Consequently, the vanishing polynomials are:
    \begin{itemize}
        \item $s_0(x) = x$
        \item $s_1(x) = x+x^2$
        \item $s_2(x) = x^4+x$
        \item $s_3(x) = x^8+x^4+x^2+x$
        \item $s_4(x) = x^{16}+x$
    \end{itemize}
    
    \end{example}


\subsubsection{A \textit{novel} polynomial basis}
    In order to perform the LCH FFT, the input polynomial has to be represented in a different basis. Given the Cantor basis $\mathbf{v} = (v_1, \cdots, v_m)$, we define the 
    \emph{novelpoly} basis with respect to $\mathbf{v}$ as $(X_1, \cdots, X_{n-1})$, defined as follows:
    \begin{equation}
        X_k(x) = \prod(s_i(x))^{b_i} \qquad \text{ where } k = \sum_{i=0}^{\log_2 n} b_i 2^i
    \end{equation}
    In other words, $X_k(x)$ is the product of all the $i$-th vanishing polynomials, where the $i$s are the set bits of $k$ in its binary representation.

\subsection{Divide et impera: the Butterflies}
    The general idea of the LCH FFT to evaluate $f(x)$ at all points of $W_k$ is to split it into $W_{k-1}$ and $W_k\setminus W_{k-1} = \{x + v_k | x \in W_{k-1}\}$:
    because of the linearity of the vanishing polynomials, the evaluation of at $W_k\setminus W_{k-1}$ can be efficiently computed by reusing the evaluation at $W_{k-1}$
    with a \emph{butterfly}.

    \begin{figure}
    \begin{tikzcd}
        {p_0(X)} &&& \oplus &&& {h_0(X)} \\
        \\
        {p_1(X)} && \bullet && \oplus && {h_1(X)}
        \arrow[from=1-1, to=1-4]
        \arrow[from=1-4, to=1-7]
        \arrow[from=3-1, to=3-3]
        \arrow["s_k(\alpha)", from=3-3, to=1-4]
        \arrow[from=3-3, to=3-5]
        \arrow[from=1-4, to=3-5]
        \arrow[from=3-5, to=3-7]
    \end{tikzcd}
    \caption{A single butterfly unit}
\end{figure}

    Similarly to other FFT algorithms, the LCH FFT is a divide and conquer algorithm: the idea is to split $f(x)$, of degree $n-1$, into two polynomials of degree $n/2-1$,
    on which we can recursively apply the same algorithm; note that the multiplication by $s_k(v_k)$ on line \ref{line:skvk} can be skipped, as it is always equal to 1.

\begin{algorithm2e}[H]
    \caption{LCH FFT}
    \KwIn{$f(x) = g(X) = g_0 + g_1X_1(x) + \cdots + g_{n-1}X_{n-1}(x)$,\newline
            $\alpha \in \mathds{F}_{2^m}$}
    \KwOut{$S = \{f(x) | x \in W_m\}$}
    \If{$\deg(f(x))$}{
        \KwRet{$[g_0]$}\;
    }
    $k \gets \max i \text{ s.t. } 2^i \leqslant \deg(f(x))$\;
    $g(X) \gets p_0(X) + X_{2^k}\cdot p_1(X)$\;
    $h_0(X) \gets p_0(X) + s_k(\alpha)\cdot p_1(X)$\;
    $h_1(X) \gets h_0(X) + s_k(v_k)\cdot p_1(X)$\;\label{line:skvk}
    \KwRet{$[FFT(h_0(X), \alpha), FFT(h_1(X), v_{k}+\alpha)]$}\;
\end{algorithm2e}




\subsubsection{Segmentation of large polynomials}
Suppose we are multiplying two binary polynomials $f(x)$ and $g(x) \in \mathds{F}_2[x]$, represented as sequences of $d$ bits:
    \begin{equation*}
        a(x) = a_0 + a_1x + \cdots + a_{d-1}x^{d-1}
    \end{equation*}
    In some cases, especially with high degree polynomials, it can be useful to split the polynomials into smaller segments, which can be easier to handle. 
    The standard \textit{Kronecker segmentation} for multiplying binary polynomials is performed as follows
    \begin{enumerate}
        \item Split the polynomials into $n = \lceil d/w \rceil$ $w$-bit blocks:
            \begin{align*}
                a(x) &= a_0 + a_1x + \cdots + a_{d-1}x^{d-1} \\
                &= (a_0 + \cdots + a_{w-1}x^{w-1}) + \cdots + (a_{d-w} + \cdots + a_{d-1}x^{w-1})x^{w(n-1)}
            \end{align*}

        \item Define the field $\mathds{F}_{2^{2w}}:=\mathds{F}_2[z]/g(z)$, with $g(z)$ irreducible polynomial of degree $2w$.
        \item Define the map $\psi: \mathds{F}_2[x] \rightarrow \mathds{F}_{2^{2w}[y]}$ such that:
        \begin{equation*}
            a^{\prime}(y) = \psi(a(x)) := a_0^{\prime} + a^{\prime}_1y + \cdots + a^{\prime}_{n-1}y^{n-1}
        \end{equation*}
        where $a^{\prime}_i = (a_{iw} + a_{iw + 1}z + \cdots + a_{iw+(w-1)}z^{w-1})$.
    
        \item Calculate the product $c^{\prime}(y) = a^{\prime}(y)\cdot b^{\prime}(y) = c_0^{\prime} + c^{\prime}_1y + \cdots \in \mathds{F}_{2^{2w}}[y]$ \label{list:mult}
        \item Map $z \mapsto x$ and $y \mapsto x^w$, collect terms and coefficients and find the result.
    \end{enumerate}
    At step~\eqref{list:mult} we employ the FFT-based method to perform the pointwise multiplication.